\section{A progress-preserving quotient of source continuations}
\label{sec:continuation-quotient}

The output program counter of a binary instruction need not be the literal
counter chosen by the original compiler. What matters is that it represents
the same future computation, with a progress condition sufficient to prevent
silent divergence. This section states that replacement precisely.

Let \(V\) be the finite set of source instruction occurrences. Every occurrence
has one of the labels
\[
  \mathrm{halt},\qquad \mathrm{inc}(r),\qquad
  \mathrm{decnz}(r),\qquad \mathrm{skip}.
\]
A break or continue occurrence is a skip with the successor obtained from its
actual path operand. Increment and skip have one successor, decrement has
successors ordered positive then zero in the distributed quotient certificate, and halt has none. The register file is
an arbitrary vector \(R\in\mathbb N^9\). The graph therefore defines a
deterministic register transition for every nonhalting configuration
\((v,R)\).

For the exact source, following only skip successors always reaches a
non-skip occurrence. Let \(e(v)\) be that occurrence, and let \(\rho(v)\) be
the number of skipped edges. The finite certificate establishes
\[
 \rho(v)=0\quad\hbox{at non-skip occurrences},\qquad
 \rho(v)=1+\rho(s(v)),\quad e(v)=e(s(v))
 \quad\hbox{at skip occurrences}.
 \tag{Q1}
\]
In this instance \(\max_v\rho(v)=4\). These are exact finite graph identities,
including absence of a skip-only cycle.

A stable partition \(\pi\) of the non-skip occurrences respects their
instruction labels, including the register operand, and satisfies
\[
 \pi(u)=\pi(v)
 \ \Longrightarrow\
 \pi(e(s_b(u)))=\pi(e(s_b(v)))
 \tag{Q2}
\]
for each corresponding successor branch \(b\). The zero and positive
successors of a decrement are compared in their specified order.
The class of an arbitrary source occurrence is
\(\bar\pi(v)=\pi(e(v))\).

\paragraph{Source equivalence.}
The quotient register machine has configurations
\((\bar\pi(v),R)\); it executes the visible increment or decrement in the
common class label, then passes to the class of the normalized successor.
A halt class is terminal. Stability makes this transition well-defined:
the labels select the same physical register and guard and induce the same
register update. Induction on a finite halting computation proves
\[
 (v,R)\text{ halts in the source}
 \quad\Longleftrightarrow\quad
 (\bar\pi(v),R)\text{ halts in the quotient}.
 \tag{Q3}
\]
The skip chains in (Q1) are finite, so neither direction discards a
divergent silent computation.

The independently reconstructed source has \(6{,}459\) internal decoder
nodes, \(6{,}460\) instruction occurrences, and \(8{,}390\) guarded successor
edges. The stable weak partition has \(4{,}530\) classes. The rank
histogram, for ranks \(0,1,2,3,4\), is respectively
\[
  4{,}533,\quad1{,}804,\quad121,\quad1,\quad1.
\]
These counts describe the source continuation certificate, not a Turing
machine state bound.

\paragraph{Binary frame relation.}
Write \(\mathcal F(v,R)\) for the set of complete canonical dispatch
configurations whose next decoded source occurrence is \(v\) and whose nine
primary registers are exactly \(R\). The frame relation includes the
complete tape, a clean path with no nonzero residual suffix beyond the
decoded occurrence, adequate separator distance, and the permitted nonempty
spare suffix. Its allowed translation and spare freedom are inherited
from the complete register-frame proof; they are not extra unconstrained
erasures.

For a nonhalting occurrence \(v\), let \(s(v,R)\) and \(R'\) be its exact
guarded source successor and updated register vector. A proposed binary
table must prove the following contract for every frame in
\(\mathcal F(v,R)\): it has a positive finite literal run to a frame in
\(\mathcal F(w,R')\), where
\[
 \bar\pi(w)=\bar\pi(s(v,R)),\qquad
 \rho(w)\leq \rho(s(v,R)).
 \tag{Q4}
\]
At a literal source halt occurrence the required contract is a finite
binary halting run. The contracts are required for every possible actual
starting occurrence, not merely one representative of each quotient class.

\begin{theorem}[Halting equivalence under quotient return contracts]
Assume the exact finite graph conditions \emph{(Q1)--(Q2)} and all complete
binary contracts \emph{(Q4)}, including the halt contracts. From any
configuration in \(\mathcal F(v,R)\), the binary machine halts if and only if
the source register machine halts from \((v,R)\).
\end{theorem}

\begin{proof}
Suppose first that the source halts. Let \(N\) be the number of nonhalting
visible instructions on the quotient computation before its halt class.
Induct on the lexicographically ordered pair \((N,\rho(v))\).

If \(v\) is a literal halt, its binary halt contract applies. If \(v\) is a
skip, its source successor has the same normalized class and registers,
while (Q1) and (Q4) give
\[
 \rho(w)\leq \rho(s(v))=\rho(v)-1.
\]
Thus the second component decreases. If \(v\) is a visible nonhalting
instruction, stability and the exact register update give a quotient
computation from the output frame with \(N-1\) visible nonhalting
instructions remaining. The first component decreases. Each binary
segment is finite, so the induction constructs a finite binary halting run.
In this particular graph the scalar measure \(5N+\rho(v)\) also strictly
decreases at every nonhalting macro.

Conversely suppose that the binary machine halts after \(t\) transitions.
Use strong induction on \(t\). A literal source halt gives the conclusion
immediately. Otherwise its contract supplies a positive segment of length
\(k\) ending in an active complete output frame. Determinism and absorption
of halting imply \(k<t\): an active configuration cannot occur at or after
the known halting time. Apply the induction hypothesis to the remaining
\(t-k\) transitions. The source beginning at the output frame halts.
By (Q3) and the class equality in (Q4), the exact source successor
\((s(v,R),R')\) also halts, and prepending the original source step proves
that \((v,R)\) halts.
\end{proof}

The rank condition is essential to the forward implication. Positive binary
duration alone would permit infinitely many returns among control counters
in the same weak class without executing the next visible instruction.
The exact rank inequality prevents that behavior. A successful table still
needs universal instruction and register proofs and its own complete
blank-tape bootstrap; the finite search fixtures are discovery evidence
only.

